\section{BÀI 2. NHIỆT PHẢN ỨNG}
\subsection{THÍ NGHIỆM}
\subsubsection{Thí nghiệm 1: Xác định nhiệt dung của nhiệt lượng kế}
\textbf{a. Tiến trình thí nghiệm}
\begin{itemize}
    \item Dùng ống đong lấy 50ml nước nguội cho vào becher rồi để thẳng nhiệt kế đo nhiệt độ $t_1$.
    \item Dùng ống đong lấy 50ml nước sôi (60-70$^\circ$C) cho vào nhiệt lượng kế rồi đo nhiệt độ $t_2$.
    \item Dùng phễu cho nhanh 50ml nước nguội vào nhiệt lượng kế rồi lắc đều cho nhiệt độ ổn định ta đo nhiệt độ $t_3$.
    \item Sử dụng công thức: $m_0c_0 = \frac{mc(t_3 - t_1) - (t_2 - t_3)}{t_2 - t_3}$
    \item Chú ý:
    \begin{itemize}
        \item m = khối lượng 50ml nước
        \item c = 1 cal/g.độ
        \item $0 < m_0c_0 < 12$
    \end{itemize}
\end{itemize}

\textbf{b. Kết quả}
\begin{itemize}
    \item Bảng số liệu:
    \begin{table}[H]
        \centering
        \begin{tabular}{|c|c|c|c|}
            \hline
            & \textbf{Lần 1} & \textbf{Lần 2} & \textbf{Lần 3} \\
            \hline
            $t_1$ & 29 & & \\
            \hline
            $t_2$ & 66 & & \\
            \hline
            $t_3$ & 48 & & \\
            \hline
            $m_0c_0$ & 2.78 & & \\
            \hline
        \end{tabular}
    \end{table}
    \item Tính giá trị $m_0c_0$:
    $$m_0c_0 = \frac{mc(t_3 - t_1) - (t_2 - t_3)}{t_2 - t_3} = \frac{50 \cdot (48 - 29) - (66 - 48)}{66 - 48} = 2.78 \text{ (cal/độ)}$$
\end{itemize}

\subsubsection{Thí nghiệm 2: Xác định hiệu ứng nhiệt của phản ứng trung hòa dung dịch HCl và NaOH}
\textbf{a. Tiến trình thí nghiệm}
\begin{itemize}
    \item Tráng buret 25ml bằng dung dịch NaOH 1M rồi lấy 25ml dung dịch NaOH 1M cho vào becher, đo nhiệt độ $t_1$.
    \item Tráng buret 25ml bằng dung dịch HCl 1M rồi lấy 25ml dung dịch HCl 1M cho vào nhiệt lượng kế, đo nhiệt độ $t_2$.
    \item Dùng phễu cho 25ml dung dịch NaOH vào nhiệt lượng kế, đo nhiệt độ $t_3$.
    \item Sử dụng công thức: $Q = (m_0c_0 + mc)(t_3 - \frac{t_1 + t_2}{2}) \text{ (cal)}$
    \item Chú ý:
    \begin{itemize}
        \item $m_{dd} = dV_{dd} \text{ (d = 1.02 g/ml)}$
        \item $c_{dd} = 1 \text{ cal/g.độ}$
        \item $\Delta H = -\frac{Q}{n}$
    \end{itemize}
\end{itemize}

\textbf{b. Kết quả}
\begin{itemize}
    \item Bảng số liệu:
    \begin{table}[H]
        \centering
        \begin{tabular}{|c|c|c|c|}
            \hline
            & \textbf{Lần 1} & \textbf{Lần 2} & \textbf{Lần 3} \\
            \hline
            $t_1$ & 29.5 & 29.5 & \\
            \hline
            $t_2$ & 30.1 & 30.5 & \\
            \hline
            $t_3$ & 35 & 35 & \\
            \hline
            $Q$ & 279.656 & 268.9 & \\
            \hline
            $Q$ trung bình & \multicolumn{3}{c|}{274.278} \\
            \hline
            $\Delta H$ & -11186.24 & -10756 & \\
            \hline
            $\Delta H$ trung bình & \multicolumn{3}{c|}{-10971.12} \\
            \hline
        \end{tabular}
    \end{table}
    \item Tính giá trị Q và $\Delta H$:
    \begin{align*}
        V_{dd} &= 25 + 25 = 50 \text{ (ml)} \\
        m &= d \times V_{dd} = 50 \times 1.02 = 51 \text{ (g)} \\
        m_0c_0 &= 2.78 \\
        Q &= (m_0c_0 + mc)\left(t_3 - \frac{t_1 + t_2}{2}\right) \\
          &= (2.78 + 51 \times 1)\left(35 - \frac{29.5 + 30.1}{2}\right) = 279.656 \text{ (cal)}
    \end{align*}
    Phương trình hóa học:
    \begin{align*}
        &\text{NaOH} + \text{HCl} \rightarrow \text{NaCl} + \text{H}_2\text{O} \\
        &0.025 \quad 0.025 \quad 0.025 \\
        \Delta H &= -\frac{Q}{n} = -\frac{279.656}{0.025} = -11186.24 \text{ (cal/mol)}
    \end{align*}
    \item Kết luận: Vì $\Delta H < 0$ nên đây là phản ứng tỏa nhiệt.
\end{itemize}

\subsubsection{Thí nghiệm 3: Xác định nhiệt hòa tan của CuSO\textsubscript{4} khan}
\textbf{a. Tiến trình thí nghiệm}
\begin{itemize}
    \item Dùng ống đong lấy 50ml nước vào nhiệt lượng kế rồi đo nhiệt độ $t_1$.
    \item Cân 4g CuSO\textsubscript{4} khan rồi cho vào nhiệt lượng kế rồi đo nhiệt độ $t_2$.
    \item Sử dụng công thức: $Q = (m_0c_0 + mc)(t_2 - t_1) \text{ (cal)}$
    \item Chú ý:
    \begin{itemize}
        \item $m_{dd} = m_{ct} + m_{dm}$
        \item $c_{dd} = 1 \text{ cal/g.độ}$
        \item $\Delta H = -\frac{Q}{n}$
    \end{itemize}
\end{itemize}

\textbf{b. Kết quả}
\begin{itemize}
    \item Bảng số liệu:
    \begin{table}[H]
        \centering
        \begin{tabular}{|c|c|c|c|}
            \hline
            & \textbf{Lần 1} & \textbf{Lần 2} & \textbf{Lần 3} \\
            \hline
            $t_1$ & 30 & 30 & \\
            \hline
            $t_2$ & 35.5 & 35 & \\
            \hline
            m (cân) & 3.7 & 3.85 & \\
            \hline
            $\Delta t$ & 5.5 & 5 & \\
            \hline
            $Q$ & 310.64 & 283.15 & \\
            \hline
            $\Delta H$ & -13433.081 & -11767.273 & \\
            \hline
            $\Delta H$ trung bình & \multicolumn{3}{c|}{-12600.177} \\
            \hline
        \end{tabular}
    \end{table}
    \item Tính giá trị Q và $\Delta H$:
    \begin{align*}
        m_{dd} &= m_{ct} + m_{dm} = 3.7 + 50 = 53.7 \\
        n_{CuSO_4} &= \frac{3.7}{160} = 0.023125 \text{ (mol)} \\
        Q &= (m_0c_0 + mc)\Delta t \\
          &= (2.78 + 53.7 \times 1) \times 5.5 = 310.64 \text{ (cal)} \\
        \Delta H &= -\frac{Q}{n} = -\frac{310.64}{0.023125} = -13433.081 \text{ (cal/mol)}
    \end{align*}
    \item Kết luận: Vì $\Delta H < 0$ nên đây là phản ứng tỏa nhiệt.
\end{itemize}

\subsubsection{Thí nghiệm 4: Xác định nhiệt hòa tan NH\textsubscript{4}Cl khan}
\textbf{a. Tiến trình thí nghiệm}
\begin{itemize}
    \item Tương tự thí nghiệm 3 nhưng thay CuSO\textsubscript{4} khan bằng NH\textsubscript{4}Cl khan.
    \item Sử dụng công thức tương tự thí nghiệm 3.
\end{itemize}

\textbf{b. Kết quả}
\begin{itemize}
    \item Bảng số liệu:
    \begin{table}[H]
        \centering
        \begin{tabular}{|c|c|c|c|}
            \hline
            & \textbf{Lần 1} & \textbf{Lần 2} & \textbf{Lần 3} \\
            \hline
            $t_1$ & 30 & 30 & \\
            \hline
            $t_2$ & 27 & 26 & \\
            \hline
            m (cân) & 3.55 & 3.59 & \\
            \hline
            $\Delta t$ & -3 & -4 & \\
            \hline
            $Q$ & -168.99 & -225.48 & \\
            \hline
            $\Delta H$ & 2546.75 & 3360.217 & \\
            \hline
            $\Delta H$ trung bình & \multicolumn{3}{c|}{2953.484} \\
            \hline
        \end{tabular}
    \end{table}
    \item Tính giá trị Q và $\Delta H$:
    \begin{align*}
        m_{dd} &= m_{ct} + m_{dm} = 3.55 + 50 = 53.55 \\
        n_{NH_4Cl} &= \frac{3.55}{53.5} = 0.066 \text{ (mol)} \\
        Q &= (m_0c_0 + mc)\Delta t \\
          &= (2.78 + 53.55 \times 1) \times -3 = -168.99 \text{ (cal)} \\
        \Delta H &= -\frac{Q}{n} = -\frac{-168.99}{0.066} = 2546.75 \text{ (cal/mol)}
    \end{align*}
    \item Kết luận: Vì $\Delta H > 0$ nên đây là phản ứng thu nhiệt.
\end{itemize}

\subsection{TRẢ LỜI CÂU HỎI}
\begin{enumerate}
    \item \textbf{$\Delta H_{tb}$ của phản ứng HCl + NaOH $\rightarrow$ NaCl + H\textsubscript{2}O sẽ được tính theo số mol HCl hay NaOH khi cho 25 ml dd HCl 2M tác dụng với 25 ml dd NaOH 1M. Tại sao?}
    \begin{itemize}
        \item $n_{HCl} = 0.025 \times 2 = 0.05 \text{ (mol)}$
        \item $n_{NaOH} = 0.025 \times 1 = 0.025 \text{ (mol)}$
        \item Ta thấy $n_{HCl} > n_{NaOH}$ nên NaOH phản ứng hết còn HCl thì còn dư $\rightarrow \Delta H_{tb}$ phản ứng được tính dựa trên số mol chất phản ứng nên sẽ được tính theo số mol NaOH.
    \end{itemize}
    \item \textbf{Nếu thay HCl 1M bằng HNO\textsubscript{3} 1M thì kết quả thí nghiệm 2 có thay đổi hay không?}
    Kết quả thí nghiệm 2 không thay đổi vì HNO\textsubscript{3} cũng là một axit mạnh, phân li hoàn toàn và tác dụng với NaOH cũng là một phản ứng trung hòa.
    \item \textbf{Tính $\Delta H_3$ bằng lý thuyết theo định luật Hess. So sánh với kết quả thí nghiệm. Hãy xem 6 nguyên nhân có thể gây ra sai số trong thí nghiệm này?}
    \begin{itemize}
        \item Mất nhiệt độ do nhiệt lượng kế.
        \item Do nhiệt kế.
        \item Do dụng cụ đong thể tích hóa chất.
        \item Do cân.
        \item Do sunfat đồng bị hút ẩm.
        \item Do lấy nhiệt dung riêng của dung dịch sunfat đồng bằng 1 cal/mol.độ.
    \end{itemize}
    \textbf{Theo em sai số nào là quan trọng nhất, giải thích? Còn nguyên nhân nào khác không?}
    \begin{itemize}
        \item Theo định luật Hess: $\Delta H_3 = \Delta H_1 + \Delta H_2 = -18.7 + 2.8 = -15.9 \text{ (kcal/mol)}$
        \item Theo thực nghiệm: $\Delta H_3 = -14.5488 \text{ (kcal/mol)}$
        \item $\rightarrow$ Chênh lệch quá lớn
        \item Theo em nguyên nhân gây ra sai số có thể là:
        \begin{itemize}
            \item Mất nhiệt độ do nhiệt lượng kế do trong quá trình thí nghiệm thao tác không chính xác, nhanh chóng dẫn đến thoát nhiệt ra môi trường bên ngoài.
            \item Sunfat đồng bị hút ẩm do trong lúc cân và đưa vào thí nghiệm thao tác không nhanh chóng.
        \end{itemize}
    \end{itemize}
\end{enumerate}
